% [VERIFY] intro traces to: problem.md O1-O5,G1,G2,Insight,A1; Book.md overview; C01,C02,C03,C05,C06,C08
\section{Introduction}
\label{sec:intro}

Deep convolutional networks have driven a sequence of breakthroughs in image
classification, where evidence suggests that network depth is of crucial importance and
that features can be enriched by the number of stacked layers. A natural question
follows: is learning better networks as easy as stacking more layers? An obstacle to
answering it is that simply increasing depth does not monotonically help. When deeper
networks are able to start converging, a \degprob is exposed: as the depth increases,
accuracy gets saturated and then degrades rapidly. Crucially, this degradation is not
caused by overfitting---adding more layers to a suitably deep model leads to higher
\emph{training} error, as observed both on CIFAR-10 with plain networks of depth 20
through 110 and on ImageNet, where a 34-layer \plainnet obtains higher top-1 validation
error than an 18-layer one despite having strictly greater capacity.

The degradation is counter-intuitive because a deeper model should be constructible from
a shallower one with no loss: copy the shallower network and append identity layers, and
the constructed deeper network has, by definition, training error no worse than its
shallower source. That such a solution exists but is not found by stochastic gradient
descent indicates that the difficulty lies in \emph{optimization} rather than in
representational capacity. The usual suspect---vanishing or exploding
gradients---can be ruled out here: batch normalization~\citep{ioffe2015batchnorm} keeps
forward-signal variances non-zero and backward gradient norms healthy, and the deep
\plainnet still attains competitive accuracy, so the solver is making progress, only
slowly. Existing remedies, including improved initialization~\citep{glorot2010xavier,he2015prelu},
batch normalization, intermediate auxiliary classifiers~\citep{lee2014dsn}, and gated
highway shortcuts~\citep{srivastava2015highway,srivastava2015trainingverydeep}, allow
networks of a few tens of layers to converge but have not demonstrated monotonic
accuracy gains at extreme depth: highway gates are data-dependent and carry parameters,
and when a gate closes the layer behaves non-residually.

We address the \degprob by introducing a deep residual learning framework. Instead of
hoping a few stacked layers directly fit a desired underlying mapping $\Hmap(\xin)$, we
explicitly let these layers fit a residual mapping $\Fmap(\xin) := \Hmap(\xin) - \xin$
and recover the original mapping as $\Fmap(\xin) + \xin$. The reformulation is motivated
by the construction argument above and by the empirical observation that layer responses
in residual nets are smaller in magnitude than in plain nets, consistent with the prior
that optimal mappings in recognition tasks sit closer to identity than to zero. The
addition is realized by a parameter-free \emph{identity shortcut} that adds the block
input to its output, so the residual variant has exactly the same number of parameters,
depth, and width as the plain counterpart---isolating the effect of residual learning
from any change in capacity. Identity shortcuts add neither parameters nor computational
cost beyond an element-wise addition, and the entire network is trained end-to-end by
ordinary stochastic gradient descent with backpropagation without modifying the
optimizer.

Our experiments on the ImageNet 2012 classification dataset and on CIFAR-10, together
with object-detection transfer experiments on PASCAL VOC and COCO, support the following
contributions:

\begin{itemize}
  \item \textbf{Diagnosis of the degradation problem.} We give controlled depth scans on
        both CIFAR-10 and ImageNet showing that deeper plain networks reach higher
        training and validation error than shallower ones under identical training, and
        argue---using batch normalization and the identity-construction
        argument---that this reflects an optimization difficulty rather than overfitting
        or vanishing gradients.
  \item \textbf{A parameter-free residual learning framework.} We reformulate stacked
        layers to fit a residual mapping with identity shortcuts, eliminating the
        degradation: at matched depth on ImageNet, the residual 34-layer net improves
        over both its plain counterpart and the shallower residual 18-layer net, and
        accuracy grows with depth up to 152 layers while the deepest model still has
        lower complexity than VGG nets.
  \item \textbf{Bottleneck designs and shortcut analysis at scale.} We show that
        three-layer bottleneck blocks make 50/101/152-layer networks practical at
        complexity comparable to a 34-layer net, that identity shortcuts suffice to fix
        degradation while projection shortcuts give only marginal gains, and that a
        six-model ensemble wins ILSVRC 2015 classification.
  \item \textbf{Generality of deep residual representations.} We demonstrate successful
        training on CIFAR-10 up to 110 layers (and a 1202-layer stress test that trains
        without optimization difficulty), and show that swapping a VGG-16 detection
        backbone for a 101-layer residual net yields large COCO and PASCAL VOC gains
        attributable solely to the learned representations.
\end{itemize}
